\input format.tex

%\usepackage{arydshln}

\begin{document}

\vspace*{3mm}
%%%
\setlength{\arrayrulewidth}{0.5pt}
\fontsize{9.3pt}{17pt}\selectfont
\color{gray2}

\begin{center}
{\bf\sanhao 健康建议}
\end{center}

%\vspace*{1mm}
\begin{spacing}{1.25}
\bigskip
\noindent
根据您本次检测结果，我们从膳食、益生菌、益生元等方面为您制定了肠道菌群干预方案，建议您遵照干预方案进行调理，以促进肠道菌群平衡。并且建议您调理3个月后复检肠道菌群，以便了解干预效果，调整干预方案，帮助您持续提升肠道健康水平。

\bigskip
\tabcolsep=6pt
\fontsize{9.3pt}{9.8pt}\selectfont
\begin{tctabularx}{tabularx={p{15.4cm}<{\centering}}}
\\[-6.769pt]
\cellcolor{topcolor} \raisebox{1pt}{\color{gray0}\xiaosihao\lantxh\bf 膳食方案} \\
\end{tctabularx}

{
\lantxh
\vspace*{-4.25mm}
\noindent
\begin{longtable}{|p{1.25cm}@{}p{2cm}|p{11.732cm}|}
\hline
\parbox[c]{\hsize}{\vskip7pt\raisebox{-1.5ex}{\includegraphics[width=.9cm]{zhushi.pdf}}\vskip7pt} & \parbox[c]{\hsize}{\centering\xiaowuhao \color{gray2}\vskip7pt主食\vskip7pt} & \parbox[c]{\hsize}{\color{gray2}\vskip7pt主食宜适量增加粗粮，如糙米饭、燕麦饭、全麦面、玉米馒头等。建议全谷物和杂豆占三分之二。\vskip7pt}\\
\hline
\parbox[c]{\hsize}{\vskip7pt\raisebox{-1.5ex}{\includegraphics[width=.9cm]{roulei.pdf}}\vskip7pt} & \parbox[c]{\hsize}{\centering\xiaowuhao \color{gray2}\vskip7pt肉类\vskip7pt} & \parbox[c]{\hsize}{\color{gray2}\vskip7pt控制动物蛋白的摄入，可选择猪瘦肉、牛肉、去皮禽肉等轮替食用。避免摄入脂肪含量高的肉类，如肥猪肉、肥牛、猪脑等。\vskip7pt}\\
\hline
\parbox[c]{\hsize}{\vskip7pt\raisebox{-1.5ex}{\includegraphics[width=.9cm]{yuxialei.pdf}}\vskip7pt} & \parbox[c]{\hsize}{\centering\xiaowuhao \color{gray2}\vskip7pt鱼虾类\vskip7pt} & \parbox[c]{\hsize}{\color{gray2}\vskip7pt宜吃适量的鱼虾类优质蛋白，建议采用清蒸、白灼的烹饪方式。\vskip7pt}\\
\hline
\parbox[c]{\hsize}{\vskip7pt\raisebox{-1.5ex}{\includegraphics[width=.9cm]{danlei.pdf}}\vskip7pt} & \parbox[c]{\hsize}{\centering\xiaowuhao \color{gray2}\vskip7pt蛋奶、大豆类\vskip7pt} & \parbox[c]{\hsize}{\color{gray2}\vskip7pt宜吃蛋白，少吃蛋黄，蛋类每天一个。每天应吃相当于鲜奶300g的奶类及奶制品，可选用无糖酸奶，每天适量多摄入豆类及豆制品，如黄豆、豆腐、豆浆等。\vskip7pt}\\
\hline
\parbox[c]{\hsize}{\vskip7pt\raisebox{-1.5ex}{\includegraphics[width=.9cm]{shucaishuiguo.pdf}}\vskip7pt} & \parbox[c]{\hsize}{\centering\xiaowuhao \color{gray2}\vskip7pt蔬菜水果\vskip7pt} & \parbox[c]{\hsize}{\color{gray2}\vskip7pt每天适量多食用新鲜水果及蔬菜，其中深色蔬菜占蔬菜总量的1{/}2，水果与蔬菜不能互相代替。多食用富含膳食纤维的果蔬，如芹菜、韭菜、火龙果。\vskip7pt}\\
\hline
\parbox[c]{\hsize}{\vskip7pt\raisebox{-1.5ex}{\includegraphics[width=.9cm]{you.pdf}}\vskip7pt} & \parbox[c]{\hsize}{\centering\xiaowuhao \color{gray2}\vskip7pt食用油、盐\vskip7pt} & \parbox[c]{\hsize}{\color{gray2}\vskip7pt食用油每天不超过25g，以橄榄油、菜籽油等植物油为宜。食盐每天不超过6g。\vskip7pt}\\
\hline
\parbox[c]{\hsize}{\vskip7pt\raisebox{-1.5ex}{\includegraphics[width=.9cm]{other.pdf}}\vskip7pt} & \parbox[c]{\hsize}{\centering\xiaowuhao \color{gray2}\vskip7pt其他\vskip7pt} & \parbox[c]{\hsize}{\color{gray2}\vskip7pt建议保持规律就餐的饮食习惯。控制总热量的摄入，保持正常体重。尽量减少外出就餐的次数。忌高脂、高热量饮食。忌暴饮暴食。戒烟限酒。如患痛风、溃疡性结肠炎、慢性腹泻、食物不耐受等有饮食禁忌的疾病，请优先遵循疾病的饮食原则。\vskip7pt}\\
\hline
\end{longtable}
\vspace*{-6mm}
\begin{longtable}{p{0.1cm}p{15.1cm}}
\includegraphics{xiaofangkuai.pdf} & {\qihao 以上膳食方案仅供参考，您的实际饮食请结合自身具体情况调整。} \\
\end{longtable}
}
\end{spacing}
\end{document}
